\documentclass[aspectratio=169]{ctexbeamer}
\usepackage[T1]{fontenc}
\usepackage{latexsym,amsmath,xcolor,multicol,booktabs,calligra}
\usepackage{graphicx,listings,stackengine}
\usepackage{hyperref}

\usepackage{PekingU}

% 去掉每节/子节前面的自动目录页
\AtBeginSection[]{}
\AtBeginSubsection[]{}

% ===== 每节课改这里 =====
\author{Cap1taL}
\title{字符串算法}
\date{2026年8月}
% ========================

% 代码高亮设置
\definecolor{deepblue}{rgb}{0,0,0.5}
\definecolor{deepred}{rgb}{0.6,0,0}
\definecolor{deepgreen}{rgb}{0,0.5,0}
\definecolor{halfgray}{gray}{0.55}

\lstset{
    basicstyle=\ttfamily\small,
    keywordstyle=\bfseries\color{deepblue},
    emphstyle=\ttfamily\color{deepred},
    stringstyle=\color{deepgreen},
    numbers=left,
    numberstyle=\small\color{halfgray},
    rulesepcolor=\color{red!20!green!20!blue!20},
    frame=shadowbox,
}

\begin{document}

\kaishu

% 封面
\begin{frame}
    \titlepage
\end{frame}

% 目录
% \begin{frame}{目录}
%     \tableofcontents[sectionstyle=show,subsectionstyle=show/shaded/hide,subsubsectionstyle=show/shaded/hide]
% \end{frame}

% ===== 从这里开始写你的内容 =====

\section{字符串算法}

\subsection{Manacher 回文串算法}
\begin{frame}{Manacher}
    \begin{itemize}
        \item 用于在线性时间内求出一个字符串内所有的回文串
        \item 回文串总量是 $O(n^2)$ 的，所以看起来在线性时间内求出所有回文串是不可能的
        \item 但事实上，我们可以在每个回文中心处记录回文半径，来将信息量降低到 $O(n)$
        \item 使用哈希加二分等手段可以在 $O(n\log n)$ 内求出它们，使用后缀数组等也可以做到线性
        \item 不过我们有针对回文串更简单的做法，就是 Manacher
        \item 下面我们将以奇数长度（回文中心是字符）的情况为例
    \end{itemize}
\end{frame}

\begin{frame}{Manacher}
    \begin{itemize}
        \item 记一个位置的回文半径为 $d_i$
        \item 我们尝试从左向右扫描，依次求出每个 $d_i$。不妨假设我们现在正要求解 $d_i$，此时所有 $<i$ 的 $d$ 都已经求解完毕了
        \item 在我们扫描的过程中，我们时刻维护我们已知的右端点 $r$ 最大的回文串
        \item 现在有多种情况：
        \begin{itemize}
            \item $i>r$：此时我们从 $[i,i]$ 开始暴力拓展，即依次考虑 $[i-1,i+1],[i-2,i+2]\cdots$
            \item $i\le r$：此时我们尝试利用以前的 $d_i$。找到 $i$ 的对应点 $j=l+(r-i)$。由于对称，$d_i$ 应该几乎与 $d_j$ 相同，除非 $d_j$ 超出了 $[l,r]$ 的范围，此时我们需要将 $d$ 截断，然后开始暴力扩展
        \end{itemize}
        \item 尝试分析 Manacher 的复杂度。每次我们暴力拓展都会使得 $r$ 增加，并且 $r$ 单调不降。因此 Manacher 的复杂度为线性
    \end{itemize}
\end{frame}

\begin{frame}{Manacher}
    \begin{itemize}
        \item 对于偶数长度的回文串
        \item 可以修改算法细节使其适用
        \item 更简单的做法是，在原串相邻字符中间插入分隔符，并直接运行奇数算法
    \end{itemize}
\end{frame}

\subsection{Trie 字典树}
\begin{frame}{01 - Trie}
    \begin{itemize}
        \item 一种特殊的 Trie，只有 $0,1$ 两种字符
        \item 01-Trie 是解决一类位值问题的利器 
    \end{itemize}
\end{frame}

\section{习题}

\subsection{洛谷 P4555 [国家集训队] 最长双回文串}

\begin{frame}{洛谷 P4555 [国家集训队] 最长双回文串}
    给定长度为 $n$ 的串 $S$，求 $S$ 的最长双回文子串 $T$，即可将 $T$ 分为两部分 $X,Y$（$|X|,|Y|\ge 1$）且 $X$ 和 $Y$ 都是回文串。
    
    $2\le n\le 10^5$
\end{frame}

\begin{frame}{Solution}
    \begin{itemize}
        \item 首先求出每个位置的回文半径，可以使用哈希或者 Manacher
        \item 考虑如何处理双回文串
        \item 我们可以枚举两个回文中心，它们可以组成一个双回文串的前提是，回文半径之和不小于间距。暴力枚举复杂度为 $O(n^2)$
        \item 注意到答案是两个回文中心间距的二倍，于是我们总是希望它们相距尽可能远
        \item 扫描线。扫描枚举右侧的回文中心 $i$，考虑左侧所有的 $[j-d_j,j+d_j]$，维护一个 $j+d_j$ 的单调栈，在单调栈上二分找到最靠左的 $j$ 即可。复杂度 $O(n\log n)$
    \end{itemize}
\end{frame}

\subsection{洛谷 P4551 最长异或路径}
\begin{frame}{洛谷 P4551 最长异或路径}
    给定一棵 $n$ 个点的带权树，求任意两个结点之间唯一路径上所有边权的异或最大值。
    
    $1\le n\le 10^5$，$0\le w<2^{31}$
\end{frame}
\begin{frame}{Solution}
    \begin{itemize}
        \item 注意到异或可以相互抵消
        \item 预处理每个点到根的异或和，问题转化为给定 $n$ 个数，要求从中找两个数，使得异或和最大
        \item 将所有数按二进制位从高到低，挂上 01 - Trie。
        \item 异或为 $1$ 当且仅当两个位不同
        \item 因此我们枚举一个数 $x$，去 01-Trie 中查询一个数与它异或和最大，只需要每步尽量走与 $x$ 位不同的边即可
        \item 复杂度 $O(n\log V)$
    \end{itemize}
\end{frame}

\subsection{洛谷 P6824 [EZEC-4] 可乐}

\begin{frame}{洛谷 P6824 [EZEC-4] 可乐}
    给定 $n$ 个正整数 $a_i$ 和整数 $k$。选择一个非负整数 $x$，最大化满足 $a_i \oplus x \le k$ 的可乐箱数。
    
    $1\le n,k,a_i\le 10^6$
\end{frame}

\begin{frame}{Solution}
    \begin{itemize}
        \item 建出 01 -Trie，在一位上异或 $ 1 $ 等价于交换左右子树
        \item 01-Trie 上 DP
        \item 设 $f_u$ 表示 $u$ 子树内，在 $k$ 的限制下能获取的最大可乐数
        \item 若 $k$ 当前位是 $ 0 $，则我们总是只能获取当前位为 $0$ 的可乐。有 $f_u=max(f_{lc_u},f_{rc_u})$
        \item 若 $k$ 当前位是 $ 1 $，则我们总是可以获取当前位为 $0$ 的可乐，而其余可乐使用子问题决策。有 $f_u=max(f_{lc_u}+sz_{rc_u},f_{rc_u}+sz_{lc_u})$
        \item 复杂度 $O(n\log V)$
    \end{itemize}
\end{frame}

% ===== 结束页 =====

\begin{frame}
    \begin{center}
        {\Huge\calligra Thanks!}
    \end{center}
\end{frame}

\end{document}
